\section{Methodology}
The overall system operates as a compilation and inference pipeline. First, an 
annotated DSL specification is compiled into a 
verification structure referred to as the Syntax Oracle. The Oracle 
enforces both syntactic and semantic constraints during text generation. Second, 
the Oracle is integrated into the inference loop of an LLM, 
constraining token generation through dynamic masking and post-hoc semantic 
validation. Finally, the Oracle is used offline to synthesize training data 
used to fine-tune the base model through LoRA~\cite{hu2021lora}, enabling 
the model to internalize structural constraints.

\subsection{Specification Language and Schema Compilation}
Rather than relying on a statically defined grammar, the proposed system 
introduces a specification protocol that extends JSON Schema with 
CSI annotations. These annotations allow 
domain constraints to be expressed directly within the schema definition (e.g., 
relational constraints such as \texttt{max\_instances $\ge$ min\_instances}).

The annotated schema is parsed by a Schema Compiler, which generates an 
executable verification structure referred to as the Syntax Oracle. The Oracle is 
formally defined as:

\begin{equation}
O = \langle M_{\text{syn}}, E_{\text{sem}} \rangle
\end{equation}

where $M_{\text{syn}}$ enforces syntactic validity and $E_{\text{sem}}$ 
evaluates semantic constraints.

\subsubsection{Syntactic Module}
The syntactic component $M_{\text{syn}}$ is constructed by translating the 
schema's CFG into a deterministic FSM:

\begin{equation}
M_{\text{syn}} = (Q, \Sigma, \delta, q_0, F)
\end{equation}

where:

\begin{itemize}
\item $Q$ is the set of parser states,
\item $\Sigma$ corresponds to the model tokenizer vocabulary $V$,
\item $\delta : Q \times \Sigma \rightarrow Q$ defines valid transitions,
\item $q_0$ is the initial state, and
\item $F$ denotes the set of terminal states.
\end{itemize}

This FSM encodes the structural constraints of the DSL and ensures that 
generated sequences remain syntactically valid, utilizing the efficient guided generation 
principles established in~\cite{willard2023efficient}.

\subsubsection{Semantic Module}
To capture cross-field dependencies without incurring exponential state 
growth, semantic constraints are represented as a directed dependency graph:

\begin{equation}
G = (V, E)
\end{equation}

where vertices $V$ correspond to generated fields and edges $E$ represent 
logical constraints between fields. Each edge $e \in E$ is compiled into 
an executable validation function.

Rather than evaluating all constraints continuously, semantic checks are 
triggered only when the FSM reaches designated terminal states. This 
event-driven design minimizes computational overhead while ensuring that 
completed DSL elements satisfy the specified invariants.

\subsection{Integration with Autoregressive Generation}
The Syntax Oracle is integrated into the inference process of the language 
model by modifying the generation loop prior to token sampling. Specifically, 
a custom logits processor---similar to those utilized in the Transformers library~\cite{wolf2020transformers}--
intercepts the model's output distribution and 
applies syntactic constraints. Let

\begin{equation}
z_t \in \mathbb{R}^{|V|}
\end{equation}

denote the raw logit vector produced by the model at timestep $t$. The 
constrained logits are computed as:

\begin{equation}
\tilde{z}_t = z_t + M_{\text{cfg},t}
\end{equation}

where $M_{\text{cfg},t} \in \{0, -\infty\}^{|V|}$ is a syntax mask 
derived from the current FSM state. Invalid tokens are assigned $-\infty$, 
preventing their selection during sampling.

\subsection{Semantic Validation and Rollback}
When the FSM reaches a designated terminal state, the generated token 
subsequence corresponding to a DSL element is extracted and evaluated 
by the semantic module. Let

\begin{equation}
Y_{\text{node}} = \{y_i, y_{i+1}, \dots, y_j\}
\end{equation}

denote the token span representing a completed element.

If the semantic constraints are satisfied, generation proceeds normally. 
Otherwise, the inference process performs a rollback by truncating the model's 
key-value cache to the state preceding the invalid sequence. The model is then 
forced to regenerate the element under additional constraints derived from the 
failed validation.

This rollback mechanism, inspired by speculative decoding techniques~\cite{leviathan2023fast}, 
ensures that generated outputs remain consistent with the specified schema while allowing 
the model to explore alternative valid continuations.

\subsection{Oracle-Guided Model Adaptation}
While runtime verification guarantees structural validity, repeated rollbacks 
can degrade generation efficiency. To mitigate this issue, the Oracle is used 
offline to synthesize a dataset of valid DSL instances.

Dataset generation is performed using coverage-guided traversal of the 
dependency graph $G$, ensuring that diverse structural configurations and 
edge cases are represented. The resulting dataset $D$ is used to fine-tune 
a pre-trained language model using Low-Rank Adaptation (LoRA)~\cite{hu2021lora}.

Let the frozen base weight matrix be

\begin{equation}
W_0 \in \mathbb{R}^{d \times k}.
\end{equation}

LoRA introduces a trainable low-rank update:

\begin{equation}
W = W_0 + \Delta W = W_0 + BA
\end{equation}

where

\begin{equation}
B \in \mathbb{R}^{d \times r}, \quad A \in \mathbb{R}^{r \times k}
\end{equation}

and $r$ is the rank of the adaptation matrices.

Training optimizes the standard causal language modeling objective:

\begin{equation}
\mathcal{L} = - \sum_{t=1}^{T} \log P_{\theta}(y_t \mid y_{<t}, X)
\end{equation}

where $X$ represents the natural language instruction and $y_t$ denotes 
tokens drawn from Oracle-verified sequences.

Through this process, the model learns to internalize common schema 
constraints, reducing the frequency of runtime corrections during inference.

\subsection{Implementation Details}
The prototype system will be implemented using a hybrid Python/C++ architecture. 
The language model and training pipeline will be developed in Python using the PyTorch~\cite{paszke2019pytorch} 
framework and the HuggingFace Transformers library. The Syntax Oracle and Schema Compiler 
will be implemented in C++ to enable efficient finite-state execution and low-latency 
constraint evaluation. Interoperability between the neural inference engine and the Oracle 
will be achieved using \texttt{pybind11}~\cite{jakob2024pybind11}, allowing direct access to tensors and enabling 
modification of model logits during generation.

Experiments will be conducted using an open-weight transformer model (e.g., LLaMA-class models~\cite{grattafiori2024llama}) 
deployed on GPU hardware. All components of the system, including schema specifications, 
generated datasets, and training configurations, will be version-controlled to 
ensure reproducibility.
